\documentclass[11pt,a4paper]{article}
\usepackage[utf8]{inputenc}
\usepackage[T1]{fontenc}
\usepackage[spanish]{babel}
\usepackage{lmodern}
\usepackage[margin=2cm]{geometry}
\usepackage{longtable,booktabs,array,xcolor,hyperref}
\usepackage{enumitem}
\hypersetup{colorlinks=true,urlcolor=blue}
\renewcommand{\arraystretch}{1.2}
\begin{document}
\begin{center}
{\Large\bfseries FORM-TL-01: Taller de tipos de requisitos}\\[5pt]
Ingeniería de Requerimientos (ISR-401) -- UTEQ\\
\textbf{Estudiante:} Mayra Lucila Córdova Carreño\\
\textbf{Docente:} Gleiston Cicerón Guerrero Ulloa, PhD\\
\textbf{Caso:} Plataforma de gestión de citas médicas -- Policlínico Ribera del Quevedo
\end{center}
\section{Alcance y criterio}
Se analizan los veinte enunciados proporcionados en la guía FORM-TL-01 [1] en tres dimensiones: tipo de requisito (funcional, calidad o restricción), producto/proceso y sistema/software. Las clasificaciones son propuestas razonadas; en casos ambiguos se explicitan los supuestos. Para requisitos de calidad se utiliza la terminología de ISO/IEC 25010:2023 [2].
\section{Clasificación y justificación}
\small
\begin{longtable}{@{}p{0.08\linewidth}p{0.12\linewidth}p{0.12\linewidth}p{0.12\linewidth}p{0.46\linewidth}@{}}
\toprule Código & Tipo & Ámbito & Nivel & Justificación (una línea)\\\midrule\endfirsthead
\toprule Código & Tipo & Ámbito & Nivel & Justificación (una línea)\\\midrule\endhead
E01 & Funcional & Producto & Software & Permite reservar citas seleccionando especialidad, médico, fecha y hora disponibles. \\
E02 & Funcional & Producto & Software & Define el envío automático de recordatorios por SMS 24 horas antes de la cita. \\
E03 & Calidad & Producto & Software & Exige un tiempo de respuesta adecuado para las operaciones de la aplicación. \\
E04 & Funcional & Producto & Software & Permite al recepcionista registrar la llegada de los pacientes. \\
E05 & Restricción & Proceso & Software & Impone Java 21 y PostgreSQL como tecnologías de desarrollo. \\
E06 & Funcional & Producto & Software & Permite consultar pacientes del día y registrar consultas atendidas. \\
E07 & Calidad & Producto & Software & Exige proteger la confidencialidad de los datos de salud. \\
E08 & Restricción & Producto & Software & Fija Android 10 o superior como plataforma mínima compatible. \\
E09 & Funcional & Producto & Software & Define el registro de inasistencia según la anticipación de cancelación. \\
E10 & Calidad & Producto & Software & Exige que los adultos mayores puedan usar la plataforma con facilidad. \\
E11 & Restricción & Proceso & Sistema & Establece seis meses como plazo máximo de entrega del proyecto. \\
E12 & Funcional & Producto & Software & Define el cierre automático de sesión tras diez minutos de inactividad. \\
E13 & Funcional & Producto & Sistema & Exige entrada por pantalla táctil e impresión de tickets en el quiosco. \\
E14 & Calidad & Producto & Software & Exige mantener disponible el servicio de la plataforma. \\
E15 & Funcional & Producto & Software & Permite generar reportes mensuales de inasistencias por especialidad. \\
E16 & Restricción & Proceso & Sistema & Obliga a validar incrementos con el director médico cada dos semanas. \\
E17 & Restricción & Producto & Sistema & Exige cumplir la normativa ecuatoriana de protección de datos personales. \\
E18 & Calidad & Producto & Sistema & Exige interoperabilidad con el sistema de facturación existente. \\
E19 & Calidad & Producto & Sistema & Exige recuperar el servicio en menos de 30 minutos sin perder citas. \\
E20 & Restricción & Proceso & Software & Obliga a revisión por pares antes de integrar cambios a la rama principal. \\
\bottomrule
\end{longtable}
\normalsize
\section{Requisitos de calidad: clasificación y cuantificación}
Los umbrales siguientes son \textbf{propuestas de especificación}, no valores aportados por el policlínico ni resultados de pruebas ejecutadas. Las métricas y pruebas deberán validarse con las partes interesadas.
\begin{enumerate}[leftmargin=*]
\item \textbf{E03 -- Eficiencia de desempeño / comportamiento temporal.} Al menos el 95\% de las consultas de disponibilidad de citas deberá responder en un máximo de 2 segundos con 300 usuarios concurrentes. \emph{Métrica:} porcentaje de consultas dentro del límite; \emph{unidades:} \% y segundos; \emph{umbral:} $\geq95\%$ en $\leq2$ s. \emph{Verificación:} prueba de carga con 300 usuarios.
\item \textbf{E07 -- Seguridad / confidencialidad.} La plataforma deberá rechazar el 100\% de los intentos de acceso a datos de salud realizados sin autorización, dentro del conjunto definido de pruebas de permisos. \emph{Métrica:} tasa de rechazos correctos; \emph{unidad:} \%; \emph{umbral:} 100\%. \emph{Verificación:} pruebas de autorización por roles. Esto comprueba solo una parte de la seguridad.
\item \textbf{E10 -- Capacidad de interacción / operabilidad.} Al menos el 90\% de los participantes mayores de 60 años deberá completar una reserva sin asistencia en un máximo de 5 minutos. \emph{Métrica:} tasa de finalización sin ayuda; \emph{unidades:} \% y minutos; \emph{umbral:} $\geq90\%$ en $\leq5$ min. \emph{Verificación:} prueba de usabilidad con participantes de esa edad, aún no realizada.
\item \textbf{E14 -- Fiabilidad / disponibilidad.} La plataforma deberá alcanzar una disponibilidad mensual de al menos el 99,5\% durante las 24 horas del día, medida sobre el servicio de consulta y reserva de citas. \emph{Métrica:} tiempo disponible sobre tiempo total; \emph{unidad:} \%; \emph{umbral:} $\geq99,5\%$. \emph{Verificación:} monitoreo mensual. La cobertura 24/7 es una propuesta que debe validarse con el director médico.
\item \textbf{E18 -- Compatibilidad / interoperabilidad.} Al menos el 99\% de los intercambios realizados con el sistema de facturación a través de su servicio web documentado deberá completarse correctamente durante las pruebas de integración. \emph{Métrica:} intercambios correctos sobre intercambios intentados; \emph{unidad:} \%; \emph{umbral:} $\geq99\%$. \emph{Verificación:} pruebas con las operaciones reales del servicio, aún por confirmar.
\end{enumerate}
\paragraph{E19 (ya medible).} Se clasifica como \textbf{Fiabilidad / recuperabilidad}: restablecimiento en \emph{menos de} 30 minutos y cero citas confirmadas perdidas tras una falla del servidor principal. Se relaciona con tolerancia a fallos, pero su foco es la recuperación.
\section{Discusión y revisión de clasificaciones}
\textbf{Estado de la comparación entre compañeros:} no se ha realizado una comparación efectiva con otro estudiante; por tanto, no se registran discrepancias ni acuerdos de revisión entre pares. Esta revisión individual \textbf{no sustituye} la comparación solicitada en la guía.

Se identifican como interpretaciones discutibles: \textbf{E05}, que puede entenderse como restricción del proceso o de implementación del producto; \textbf{E08}, como restricción de plataforma o característica de adaptabilidad; \textbf{E11}, cuyo nivel sistema/software no se deduce directamente del plazo del proyecto; y \textbf{E16}, cuya validación puede abarcar todo el sistema o solo incrementos de software. Se conservan las clasificaciones de la tabla como propuestas y no se atribuyen a ningún compañero.
\section{Uso de inteligencia artificial}
Se utilizó \textbf{ChatGPT (OpenAI, GPT-6)} como apoyo para estructurar, redactar y revisar propuestas de clasificación, métricas y texto LaTeX. La estudiante debe revisar el contenido contra la guía del docente, confirmar los supuestos y verificar la compilación y las fuentes antes de entregarlo. No se realizaron ni se afirman pruebas de campo o discusiones entre pares.
\section{Referencias}
\begin{enumerate}[label={[\arabic*]}]
\item G. C. Guerrero-Ulloa, ``Guía y rúbrica de evaluación: FORM-TL-01, Taller de tipos de requisitos,'' Universidad Técnica Estatal de Quevedo, material de clase, 2026.
\item ISO/IEC, \emph{ISO/IEC 25010:2023, Systems and software engineering -- Systems and software Quality Requirements and Evaluation (SQuaRE) -- Product quality model}, 2023.
\end{enumerate}
\end{document}
